\section{Product distances and measures across scales}
\label{sec:scales}

Suppose, toward a contradiction, that there is a sequence of parameters
$k\to+\infty$ for which the top Lyapunov exponent vanishes almost
everywhere. All limits in this section are taken along this sequence,
passing to subsequences when indicated, except where a fixed-parameter
limit in orbit length is explicitly specified. The structural constants
in \cref{can:theorem,bridge:theorem} are fixed throughout; the bridge
estimate's auxiliary constants may depend on its observation bounds.
Use the growth semimetric $g$ and expected shortfall $e_n$ from
\cref{scale:distance-deficit}; their parameter dependence remains
suppressed. At each parameter in the sequence, the zero-exponent
assumption and bounded convergence give
\begin{equation}\label{scale:long-deficit}
e_n\longrightarrow1\qquad(n\to\infty).
\end{equation}
By contrast, the shortfall at every fixed length tends to zero as the
parameter grows. We will select a dyadic range in which the shortfall
first increases, and sum the loss between each whole interval and its
two halves over that range. This sum telescopes, but its limiting mass
must still be controlled near arrays with one constant growth rate.
The local tests and measures below are designed to retain that mass.

\Needspace{8\baselineskip}
\subsection{Four-point comparison and truncation}

\begin{lemma}[Four-point comparison]\label{scale:fourpoint}
Fix a base index $i$ and write
\[
 D_j=\norm{T_{i,j}},\qquad
 R_{jp}=M^{g(j,p)-g(i,j)-g(i,p)}.
\]
There is an absolute constant $C_0\ge1$ such that, for every orbit,
every sufficiently large parameter, and every choice of integer indices,
\begin{equation}\label{scale:R-fourpoint}
 R_{jp}\le C_0\max\{R_{jl},R_{lp}\}.
\end{equation}
\end{lemma}

\begin{proof}
Choose a most expanded unit input direction $a_j$ for $T_{i,j}$;
if $D_j=1$, any singular direction may be used. Set
\[
 Q_{jp}=\max\{|a_j\wedge a_p|,D_j^{-2},D_p^{-2}\}.
\]
We claim
\begin{equation}\label{scale:R-comparison}
 Q_{jp}\le R_{jp}\le3Q_{jp}.
\end{equation}
To obtain the sine lower bound, use a unit input at the base
perpendicular to $a_p$. Its norm at $p$ is $D_p^{-1}$, whereas its
norm at $j$ is at least $D_j|a_j\wedge a_p|$. The transfer from $p$
to $j$ therefore has norm at least $D_jD_p|a_j\wedge a_p|$.
The inequalities
\[
 \norm{T_{p,j}}\ge D_j/D_p,\qquad
 \norm{T_{p,j}}\ge D_p/D_j
\]
give the other two lower bounds; the second uses equality of the norms
of a determinant-one matrix and its inverse.

For the upper bound, transport a unit vector from $p$ back to the
base. Its components along $a_p$ and $a_p^\perp$ have absolute values
at most $D_p^{-1}$ and $D_p$, respectively. Also
\[
 \norm{T_{i,j}a_p^\perp}
 \le D_j|a_j\wedge a_p|+D_j^{-1}.
\]
It follows that
\[
 R_{jp}\le D_p^{-2}+|a_j\wedge a_p|+D_j^{-2},
\]
proving \cref{scale:R-comparison}. By
\cref{setup:sine-triangle},
$Q_{jp}\le2\max\{Q_{jl},Q_{lp}\}$, so $C_0=6$ is permissible.
\end{proof}

\begin{lemma}[Truncation of a cancellation]\label{scale:truncation}
Fix integers $i$, $n\ge1$, and $1\le s\le n$, and put
\[
 L=i-n,\quad R=i+n,\quad L'=i-s,\quad R'=i+s,
\]
\[
 D=n-\min\{g(i,L),g(i,R)\},\qquad
 \ell=g(i,L)+g(i,R)-g(L,R).
\]
Use the quantities $R_{uv}$ of \cref{scale:fourpoint} with base $i$.
Then
\begin{align}
 g(i,L'),g(i,R')&\ge s-D,
 \label{scale:cut-wings}\\
 R_{L'R'}&\le C_0^2\max\{R_{LR},M^{-2s+2D}\}.
 \label{scale:cut-R}
\end{align}
There is a constant $A$ independent of the parameter and indices such
that the cut cancellation loss satisfies
\begin{equation}\label{scale:cut-loss}
 g(i,L')+g(i,R')-g(L',R')
 \ge\min\{\ell,2s-2D\}-A,
\end{equation}
and consequently
\begin{equation}\label{scale:cut-full}
 g(L',R')\le2s-\min\{\ell,2s-2D\}+A.
\end{equation}
\end{lemma}

\begin{proof}
The triangle inequality and \cref{scale:distance-bound} give
\[
 g(i,L')\ge g(i,L)-g(L,L')\ge(n-D)-(n-s)=s-D,
\]
and similarly on the right. In particular,
\[
 R_{LL'}\le M^{(n-s)-(n-D)-(s-D)}=M^{-2s+2D},
 \qquad R_{RR'}\le M^{-2s+2D}.
\]
Apply \cref{scale:R-fourpoint} twice along the chain $L',L,R,R'$ to
obtain \cref{scale:cut-R}. Taking negative base-$M$ logarithms gives
\cref{scale:cut-loss}, with any fixed upper bound for
$2\log_M C_0$ as $A$. Such a bound is uniform after imposing, for
example, $M\ge2$. Finally, each cut wing has length at most $s$,
which yields \cref{scale:cut-full}.
\end{proof}

\begin{lemma}[Doubling bound]\label{scale:doubling}
For all sufficiently large parameters and every integer $n\ge1$,
\begin{equation}\label{scale:doubling-equation}
 e_n\le e_{2n}\le C e_n+\frac{C}{n},
\end{equation}
with a constant independent of $n$ and the parameter.
\end{lemma}

\begin{proof}
Use the notation of \cref{scale:truncation}. Both $D$ and $\ell$ are
nonnegative, and $\ell\le2n$. Let $\delta>0$ be the fixed constant
in \cref{can:theorem}. Choose a fixed $H$ sufficiently large that
\[
 \ell>H(1+D),\qquad s=\lfloor\ell/2\rfloor
 \quad\Longrightarrow\quad
 s\ge1,\quad 2D+A\le\delta s.
\]
This is possible since $s\ge\ell/2-1$. Also $s\le n$ and
$2s\le\ell$. Thus \cref{scale:cut-wings,scale:cut-full} imply
\[
 \min\{g(i,i-s),g(i,i+s)\}\ge(1-\delta)s,
 \qquad g(i-s,i+s)\le2D+A\le\delta s.
\]
This is the event bounded by \cref{can:theorem}, with spacing $s$.
On the part where this integer $s$ is selected, $\ell<2s+2$.
Writing the constants of that theorem as $C_{\rm can},\gamma>0$,
we obtain
\[
 \E\bigl[\ell\,\one_{\{\ell>H(1+D)\}}\bigr]
 \le C_{\rm can}\sum_{s=1}^{\infty}(2s+2)e^{-\gamma s}
 \le C.
\]
The complementary contribution is at most $H(1+\E D)$, hence
\begin{equation}\label{scale:mean-cancellation}
 \E\ell\le C(1+\E D).
\end{equation}
Stationarity gives
\[
 \E D\le\E[n-g(i,L)]+\E[n-g(i,R)]=2ne_n,
 \qquad e_{2n}=e_n+\frac{\E\ell}{2n}.
\]
These identities and \cref{scale:mean-cancellation} prove the result.
In particular the lower inequality uses only $\ell\ge0$.
\end{proof}

\Needspace{8\baselineskip}
\subsection{Selecting the first growing deficit}

\begin{lemma}[Critical scales]\label{scale:selection}
Along the assumed zero-exponent parameter sequence, there are powers
of two $n_0=n_0(K)$ and $N=N(K)>n_0$, and numbers $\epsilon=e_N$,
such that
\begin{equation}\label{scale:selected-limits}
 n_0\to\infty,\qquad \epsilon\to0,\qquad
 \epsilon n_0\to\infty,
\end{equation}
and, for every fixed integer $m\in\Z$ (once the lengths below are
positive integers),
\begin{equation}\label{scale:selected-boundaries}
 e_{2^m n_0}=o(\epsilon),\qquad
 e_{2^m N}=O_m(\epsilon).
\end{equation}
All scale sums below run over powers of two.
\end{lemma}

\begin{proof}
First, for every fixed integer $a\ge1$, we have $e_a\to0$ as
$K\to\infty$. To see this, fix $0<\zeta<1$. The uniform marginal
law of each $q_i$ implies
\[
 \PP\{|v(q_i)|<M^{1-\zeta}\}\longrightarrow0.
\]
Indeed $v(x)=2+K\cos(2\pi x)$, and the length of the set on which
$|\cos(2\pi x)|\le(M^{1-\zeta}+2)/K$ tends to zero. A union bound
therefore shows that, with probability tending to one, all $a$
coefficients in the product satisfy $|v(q_i)|\ge M^{1-\zeta}$.
Start the linear recurrence with $(u_0,u_{-1})=(1,0)$. On this event,
for $M$ large, induction gives
\[
 |u_{j+1}|\ge(M^{1-\zeta}-1)|u_j|\ge|u_j|
 \quad(0\le j<a).
\]
Consequently $g(0,a)/a\ge\log_M(M^{1-\zeta}-1)$ there.
The deficit always lies in $[0,1]$, so
$\limsup_{K\to\infty}e_a\le\zeta$. Letting $\zeta\downarrow0$
proves the claim.

Choose powers of two $n_0\to\infty$ sufficiently slowly that
$e_{n_0}\to0$, by diagonal selection over the fixed lengths
$1,2,4,\ldots$. Set
\[
 \theta_K=(n_0^{-1}+e_{n_0})^{1/2}.
\]
For all sufficiently large $K$, $e_{n_0}<\theta_K<1$.
By \cref{scale:long-deficit}, there is a first power of two
$N\ge n_0$ for which $e_N>\theta_K$. It satisfies $N>n_0$ and
$e_{N/2}\le\theta_K$. Applying \cref{scale:doubling} once gives
\[
 \theta_K<\epsilon=e_N\le C\theta_K+C/n_0\longrightarrow0.
\]
Also $\epsilon n_0\ge\theta_K n_0\ge\sqrt{n_0}\to\infty$.
For fixed $m\ge0$, iteration of \cref{scale:doubling} yields
\[
 e_{2^m n_0}\le C_m(e_{n_0}+n_0^{-1})
 =C_m\theta_K^2=o(\epsilon),
\]
and
\[
 e_{2^m N}\le C_m(\epsilon+N^{-1})=O_m(\epsilon).
\]
For a fixed negative $m$, the lengths $2^m n_0$ and $2^mN$ are
eventually positive integers. Dyadic monotonicity bounds their
deficits by $e_{n_0}=o(\epsilon)$ and $e_N=\epsilon$, respectively.
This proves every assertion.
\end{proof}

\subsection{A compact space of distance arrays}

The next construction records shortfall increments as measures on
rescaled distance arrays. Its output will be locally finite limits
away from affine arrays, together with a dilation balance; the
discarded affine part requires separate control in the final argument.

Let $\mathcal D=\Z[1/2]$ denote the dyadic rational numbers. Define
$X$ to be the set of symmetric semimetric arrays
$d:\mathcal D\times\mathcal D\to[0,\infty)$ satisfying
\begin{equation}\label{scale:ambient-bound}
 d(s,t)\le|s-t|+2.
\end{equation}
Give $X$ the topology of pointwise convergence. It is a closed subset
of a countable product of compact intervals, and hence is compact
and metrizable. Its compact subset of affine arrays is
\[
 \mathcal A=\{d_w:0\le w\le1\},\qquad d_w(s,t)=w|t-s|.
\]
These are precisely the arrays of this form that belong to $X$:
the bound at arbitrarily large separations forces $w\le1$.
Write $\mathcal U=X\setminus\mathcal A$.

For a positive integer $n$, define the random array
\begin{equation}\label{scale:rescaled-array}
 g_n(s,t)=\frac{g(\lfloor ns\rfloor,\lfloor nt\rfloor)}{n}
 \qquad(s,t\in\mathcal D).
\end{equation}
It belongs to $X$, since its distances are at most
$|s-t|+1/n$. Introduce the finite measures
\begin{equation}\label{scale:prelimit-measures}
 \mu_K=\epsilon^{-1}\sum_{\substack{n_0\le n<N\\n\text{ a power of }2}}
                  \Law(g_n),\qquad
 \nu_K=\epsilon^{-1}\Law(g_N),\qquad
 \xi_K=\epsilon^{-1}\Law(g_{n_0}).
\end{equation}
Their total masses need not remain bounded. We will obtain limits on
$\mathcal U$ by controlling local departures from affinity.

For an interval $I=[s,t]$ with distinct dyadic endpoints, let $I_l,I_r$
be its left and right halves and define continuous functions on $X$ by
\begin{equation}\label{scale:tests}
 h(I)=1-\frac{d(s,t)}{t-s},\qquad
 J(I)=h(I)-\frac{h(I_l)+h(I_r)}{2},\qquad
 V(I)=\bigl(h(I_l)-h(I_r)\bigr)^2.
\end{equation}
The functions $J(I)$ and $V(I)$ are nonnegative on all of $X$.
The first assertion is the triangle inequality. The shortfall $h(I)$
itself need not be nonnegative on all of $X$. It does belong to $[0,1]$
on a sample $g_n$ whenever the endpoints multiplied by $n$ are integers.
For any fixed dyadic interval, this alignment holds at every scale
$n\ge n_0$ occurring here once $K$ is sufficiently large.

\begin{lemma}[Bounds for local tests]\label{scale:local-mass}
Let $I$ have dyadic endpoints and length $p=2^m$, where $m\in\Z$ is
fixed. Then
\begin{equation}\label{scale:sum-test-bound}
 \int\bigl(J(I)+V(I)\bigr)\,\dd\mu_K=O_I(1).
\end{equation}
Uniformly over the individual dyadic scales $n_0\le n\le N$,
\begin{equation}\label{scale:individual-test-bound}
 \epsilon^{-1}\E\bigl[h(I)(g_n)+J(I)(g_n)+V(I)(g_n)\bigr]
 =O_I(1).
\end{equation}
At the initial scale the stronger bound holds:
\begin{equation}\label{scale:initial-test-bound}
 \epsilon^{-1}\E\bigl[h(I)(g_{n_0})+J(I)(g_{n_0})
                         +V(I)(g_{n_0})\bigr]=o(1).
\end{equation}
All three statements concern sufficiently large $K$, so that the
indicated endpoints and half-endpoints align.
\end{lemma}

\begin{proof}
By stationarity,
\[
 \E h(I)(g_n)=e_{pn},\qquad
 \E J(I)(g_n)=e_{pn}-e_{pn/2}.
\]
Summation over consecutive dyadic scales gives the exact identity
\begin{equation}\label{scale:J-telescope}
 \sum_{\substack{n_0\le n<N\\n\text{ a power of }2}}
       \E J(I)(g_n)=e_{pN/2}-e_{pn_0/2}.
\end{equation}
By \cref{scale:selection} and dyadic monotonicity, its nonnegative
right side is $O_I(\epsilon)$.

Put $a=h(I_l)$, $b=h(I_r)$, and $\bar h=(a+b)/2$. At the aligned
sampled arrays, $0\le\bar h\le h(I)\le1$, and therefore
\begin{align*}
 \frac{V(I)}4
 &=\frac{a^2+b^2}{2}-\bar h^2\\
 &\le\frac{a^2+b^2}{2}-h(I)^2+2J(I).
\end{align*}
Write $b_q=\E\bigl[(1-g(0,q)/q)^2\bigr]$ for integer $q\ge1$.
By stationarity and another dyadic telescoping,
\[
 \frac14\sum_n\E V(I)(g_n)
 \le b_{pn_0/2}-b_{pN/2}
      +2\bigl(e_{pN/2}-e_{pn_0/2}\bigr)
 \le b_{pn_0/2}+2e_{pN/2}.
\]
Here and in this display only, the unqualified sum is over the scales
in \cref{scale:J-telescope}. Since $0\le b_q\le e_q$, the last
expression is $O_I(\epsilon)$. This proves
\cref{scale:sum-test-bound}.

At any single aligned scale, $J(I)\le h(I)$ and
$V(I)=(a-b)^2\le a+b$. Hence the expectation in
\cref{scale:individual-test-bound} is at most
$2e_{pn}+2e_{pn/2}$. Dyadic monotonicity bounds these terms by the
corresponding terminal-scale deficits, which are $O_I(\epsilon)$
by \cref{scale:selection}. At $n=n_0$ they are $o(\epsilon)$,
again by that lemma, including its consequence for negative shifts.
\end{proof}

\begin{lemma}[Detection of affinity]\label{scale:affinity-tests}
As $I$ ranges over intervals with dyadic endpoints and power-of-two
length, the common zero set of $J(I)+V(I)$ is exactly $\mathcal A$.
Consequently, every compact subset $F\subset\mathcal U$ admits
finitely many such intervals $I_1,\ldots,I_b$ and a number $c_F>0$
such that
\begin{equation}\label{scale:finite-test-cover}
 \sum_{a=1}^{b}\bigl(J(I_a)+V(I_a)\bigr)\ge c_F
 \quad\hbox{on }F.
\end{equation}
\end{lemma}

\begin{proof}
Every affine array annihilates these tests. Conversely, suppose they
all vanish. For one of the stated intervals $I$, the equalities
$J(I)=V(I)=0$ say that its two halves have the same speed as $I$
and that their distances add to the full distance. Repeating this
argument shows that every edge of every dyadic subdivision of $I$
has distance $w_I$ times its length, with
$w_I=d(s,t)/(t-s)$. The triangle inequality gives the same upper
bound between any two subdivision nodes. Applying the triangle
inequality across the full interval, using those nodes as intermediate
points, gives the reverse bound, and therefore equality. Every dyadic
point of $I$ is a subdivision node at some level. Apply this argument
to $[-2^a,2^a]$ for arbitrarily large integers $a$. The common pair
$0,1$ fixes one speed on all these intervals, so the array is globally
affine. The ambient bound forces that speed to lie in $[0,1]$.

For the last assertion, each point of $F$ has a positive test. The
corresponding open sets cover $F$, and a finite subcover suffices.
Their sum has a positive minimum on the compact set $F$.
\end{proof}

\Needspace{8\baselineskip}
\subsection{Locally finite limits and their symmetries}

The limiting inequality \cref{scale:tree-inequality} is the
zero-hyperbolic Gromov-product condition: after identifying points at
zero semimetric distance, $B_v(s,t)/2$ is the Gromov product based at $v$;
see \cite[Section~2.2, Definition~2.10]{Bestvina1997}.
The determinant-one estimate and the passage to the limiting arrays
are proved here; no real-tree representation theorem is used.

\begin{proposition}[Limits away from affinity]\label{scale:limits}
After passing to a subsequence, the restrictions of $\mu_K$ and
$\nu_K$ to $\mathcal U$ converge vaguely to locally finite
nonnegative Borel measures $\mu$ and $\nu$, respectively. Moreover,
\begin{equation}\label{scale:initial-negligible}
 \xi_K\longrightarrow0\quad\hbox{vaguely on }\mathcal U.
\end{equation}
The measures $\mu$ and $\nu$ are supported on the closed class
$X_0\subset X$ consisting of arrays satisfying $d(s,t)\le|s-t|$
and the inequalities
\begin{equation}\label{scale:tree-inequality}
 B_v(s,t)\ge\min\{B_v(s,p),B_v(p,t)\},\qquad
 B_v(s,t)=d(v,s)+d(v,t)-d(s,t)
\end{equation}
for every $v,s,t,p\in\mathcal D$. Every member of $X_0$ extends
uniquely to a continuous semimetric on $\R$, obeying the same
inequalities at all real arguments. Both limiting measures are
invariant under
\[
 (\tau_r d)(s,t)=d(s+r,t+r)\qquad(r\in\mathcal D).
\]
For every interval $I$ as in \cref{scale:local-mass},
\begin{equation}\label{scale:limiting-test-bound}
 \int(J(I)+V(I))\,\dd\mu<\infty,
 \qquad
 \int(h(I)+J(I)+V(I))\,\dd\nu<\infty.
\end{equation}
In particular,
\begin{equation}\label{scale:terminal-shortfall}
 \int h([0,1])\,\dd\nu\le1.
\end{equation}
\end{proposition}

\begin{proof}
For each compact $F\subset\mathcal U$, combine
\cref{scale:finite-test-cover} with
\cref{scale:sum-test-bound,scale:individual-test-bound} to obtain
uniform bounds on $\mu_K(F)$ and $\nu_K(F)$. Using
\cref{scale:initial-test-bound} instead gives $\xi_K(F)=o(1)$.
The latter proves \cref{scale:initial-negligible}, since a continuous
compactly supported test is bounded and has compact support.

Here is an explicit way to extract the local limits. Fix a compatible
metric on $X$ and let $a(d)=\dist(d,\mathcal A)$. Choose a continuous
nondecreasing function $\chi:[0,\infty)\to[0,1]$ equal to zero on
$[0,1/2]$ and to one on $[1,\infty)$. Set
$\chi_j(d)=\chi(2^j a(d))$. These functions increase to one on
$\mathcal U$, have compact support there, and satisfy
$\chi_{j+1}=1$ on $\supp\chi_j$. The finite measures
$\chi_j\mu_K$ and $\chi_j\nu_K$ have bounded total masses for
each $j$. Weak compactness for finite measures on compact metric
spaces, followed by a diagonal extraction, gives weak limits for
all of them. Their restrictions agree on the open regions where the
corresponding cutoffs are one: this follows by testing against
continuous functions supported in such a region, for which the
prelimit integrals agree. These compatible restrictions define
locally finite measures on $\mathcal U$. Every compactly supported
continuous function is supported where some $\chi_j$ is one, so
these are the asserted vague limits.

For completeness, the support restrictions survive despite the
possibly unbounded total masses. The arrays $g_n$ satisfy
$g_n(s,t)\le|s-t|+1/n$. Taking logarithms in
\cref{scale:R-fourpoint}, then dividing by $n$, shows that their
version of \cref{scale:tree-inequality} has error at most
$(\log_M C_0)/n$. These errors are bounded uniformly by $C/n_0$
at every sampled scale. Any fixed strict violation, by a fixed
positive margin, is thus absent from every sampled array once $K$
is large. Testing on compact subsets of its open violation set
shows that it has zero limiting measure. There are countably many
dyadic arguments and positive rational margins, which proves the
stated support assertion. For an array with $d(s,t)\le|s-t|$,
the semimetric inequality gives
\[
 |d(s,t)-d(s',t')|\le|s-s'|+|t-t'|.
\]
It therefore has a unique jointly continuous extension to real
arguments; all its inequalities extend by continuity.

For fixed $r\in\mathcal D$, each prelimit law is invariant under
$\tau_r$ once $n_0r$ is an integer. Indeed
$\lfloor n(s+r)\rfloor=\lfloor ns\rfloor+nr$, and stationarity
applies simultaneously to the array. The map $\tau_r$ is a
homeomorphism of $X$ preserving $\mathcal A$, so it preserves
compact support in $\mathcal U$. Invariance therefore passes to
the vague limits.

Finally, multiply any nonnegative test $J(I)+V(I)$ by $\chi_j$,
pass to the vague limit, and use
\cref{scale:sum-test-bound}. Monotone convergence as $j\to\infty$
gives the first bound in \cref{scale:limiting-test-bound}. For $\nu$
use \cref{scale:individual-test-bound} in the same way. Although
$h(I)$ is not nonnegative everywhere on $X$, it is nonnegative on
the aligned samples and on the supported limiting class $X_0$;
these facts justify the cutoff bounds and monotone convergence for
that term as well. In the case $I=[0,1]$ the uncut terminal
integral is exactly $e_N/\epsilon=1$, giving
\cref{scale:terminal-shortfall}.
\end{proof}

Define the dilation
\begin{equation}\label{scale:dilation}
 (Sd)(s,t)=\frac{d(2s,2t)}2.
\end{equation}
It is a continuous map from $X$ into $X$ and fixes every affine
array. Its inverse need not map all of $X$ into $X$, but on $X_0$
it is the homeomorphism
\[
 (S^{-1}d)(s,t)=2d(s/2,t/2).
\]
Both maps preserve $X_0\cap\mathcal U$.

\begin{proposition}[Positive dilation decomposition]\label{scale:decomposition}
On $\mathcal U$, with the measures supported on $X_0$ as above,
\begin{equation}\label{scale:measure-balance}
 S_*\mu=\mu+\nu.
\end{equation}
There is a locally finite nonnegative measure $\mu_\infty$, invariant
under $S$ and every dyadic translation, such that
\begin{equation}\label{scale:measure-decomposition}
 \mu=\mu_\infty+\sum_{j=1}^{\infty}(S^{-j})_*\nu,
 \qquad S_*\mu_\infty=\mu_\infty.
\end{equation}
The series is a series of positive measures dominated by $\mu$.
\end{proposition}

\begin{proof}
The floor convention in \cref{scale:rescaled-array} gives the exact
identity $Sg_n=g_{2n}$. Hence, as finite measures on $X$,
\begin{equation}\label{scale:finite-measure-balance}
 S_*\mu_K=\mu_K+\nu_K-\xi_K.
\end{equation}
To pass to the local limits, observe that $S^{-1}(\mathcal A)
=\mathcal A$: multiplication by two permutes the dyadic rationals,
so $Sd=d_w$ implies $d=d_w$. If $F\subset\mathcal U$ is compact,
its preimage under $S$ is closed in compact $X$ and disjoint from
$\mathcal A$, hence compact in $\mathcal U$. Thus $S$ is proper
on $\mathcal U$, and composition with $S$ preserves continuous
compactly supported tests there. Vague convergence and
\cref{scale:initial-negligible} now give
\cref{scale:measure-balance}.

We may henceforth work on $X_0\cap\mathcal U$, where $S$ is
invertible. Applying its inverse to \cref{scale:measure-balance}
and iterating gives, for every $b\ge1$,
\[
 \mu=(S^{-b})_*\mu+\sum_{j=1}^{b}(S^{-j})_*\nu.
\]
The increasing positive partial sums are dominated by $\mu$.
Their limit
$\lambda=\sum_{j\ge1}(S^{-j})_*\nu$ is therefore a locally finite
measure. On relatively compact Borel sets define
$\mu_\infty=\mu-\lambda$. The definitions agree on overlaps and
extend by exhaustion to a locally finite nonnegative measure;
this construction avoids subtracting infinite total masses.
Monotone convergence for the positive series yields
\[
 S_*\lambda=\nu+\lambda.
\]
Subtract this identity from \cref{scale:measure-balance} on
relatively compact sets to obtain $S_*\mu_\infty=\mu_\infty$.

For $r\in\mathcal D$ and $j\ge1$, the identity
\[
 \tau_r S^{-j}=S^{-j}\tau_{r/2^j}
\]
and translation invariance of $\nu$ show that every summand of
$\lambda$ is invariant under $\tau_r$. The same holds for their
sum and, by local subtraction from the translation-invariant
measure $\mu$, for $\mu_\infty$. This proves all assertions.
\end{proof}
